\documentclass[11pt]{article}
\usepackage[margin=1in]{geometry}
\usepackage{amsmath,amssymb,amsthm,graphicx,microtype}
\usepackage[hidelinks]{hyperref}

\newtheorem{theorem}{Theorem}
\newtheorem{corollary}[theorem]{Corollary}
\theoremstyle{remark}
\newtheorem{remark}[theorem]{Remark}
\newcommand{\N}{\mathbb N}
\newcommand{\cV}{\overline V}
\newcommand{\kp}{k_p(V)}

\title{An exact projective-weight criterion for\\
topological ergodicity of coechelon shifts}
\author{Candidate full answer to Open Problem (1) of arXiv:1911.07513}
\date{13 August 2026}

\begin{document}
\maketitle

\begin{center}
\fbox{\parbox{0.91\textwidth}{\textbf{Status.}
Candidate full solution, likely valid, of Open Problem (1).  The theorem
below gives a necessary and sufficient condition directly in the defining
weights.  It does not claim to solve the source's separate Open Problem
(2).}}
\end{center}

\section{Source question}

The source is J.~Bonet, T.~Kalmes, and A.~Peris, \emph{Dynamics of shift
operators on non-metrizable sequence spaces}, arXiv:1911.07513v2
\cite{BKP}.  It appeared in \emph{Revista Matem\'atica Iberoamericana}
37(6) (2021), 2373--2397.  Its final subsection asks for a condition that
completely characterizes topological ergodicity of the backward shift on a
K\"othe coechelon space.  Figure~\ref{fig:question} is a crop of page 20 of
the source PDF.

\begin{figure}[ht]
\centering
\includegraphics[width=\textwidth]{figures/open_problems_crop.png}
\caption{The two open problems stated in the source.  This packet answers
Problem (1).}
\label{fig:question}
\end{figure}

\section{Setting}

Let $V=(v^{(m)})_{m\in\N}$ be a decreasing sequence of strictly positive
weights on $\N$, and let $p=0$ or $1\leq p<\infty$.  Assume that the
backward shift
\[
 Be_1=0,\qquad Be_j=e_{j-1}\quad(j\geq2)
\]
is continuous on the K\"othe coechelon space $\kp$.  As in the source, its
projective family is
\[
 \cV=\left\{\bar v\in[0,\infty)^{\N}:
 \text{for every }m\text{ there is }a_m>0\text{ such that }
 \bar v_j\leq a_m v^{(m)}_j\text{ for all }j\right\}.
 \tag{2.1}
\]
The topology of $\kp=K_p(\cV)$ is generated by
\[
 q_{\bar v}(x)=
 \begin{cases}
  \bigl(\sum_j |x_j|^p\bar v_j^p\bigr)^{1/p},&1\leq p<\infty,\\[2mm]
  \sup_j |x_j|\bar v_j,&p=0,
 \end{cases}
 \qquad \bar v\in\cV.
 \tag{2.2}
\]
A set $A\subseteq\N$ is \emph{syndetic} if its gaps are bounded.  The
operator $B$ is topologically ergodic when
\[
 N_B(U,W)=\{n\geq0:B^n(U)\cap W\ne\varnothing\}
 \tag{2.3}
\]
is syndetic for every pair of nonempty open sets $U,W$.

\section{Proof intuition}

For necessity, test ergodicity with a small neighborhood of zero and a
small neighborhood of $e_1$.  A return at time $n$ forces the $(n+1)$st
coordinate of the starting vector to be close to one.  Smallness in the
projective seminorm then forces the projective weight at $n+1$ to be small.
Thus the small-coordinate set of every projective weight must be syndetic.

For sufficiency, start with finitely supported points in two arbitrary open
sets.  Only finitely many inverse images of one zero neighborhood are
needed to implant the second point far out in the first one.  A finite
intersection of projective seminorm balls is dominated by one projective
seminorm ball.  The assumed syndeticity therefore supplies syndetically
many locations at which the implantation works.

Finally, all projective weights can be tested at once through the canonical
weights
\[
 \bar v^{,a}_j=\inf_m a_m v^{(m)}_j.
 \tag{3.1}
\]
This converts the topological condition into a formula involving only the
given matrix $V$.

\section{Exact characterization}

\begin{theorem}[answer to Open Problem (1)]\label{thm:main}
Under the assumptions above, the following are equivalent.
\begin{enumerate}
\item[(i)] The backward shift $B$ is topologically ergodic on $\kp$.
\item[(ii)] For every $\bar v\in\cV$ and every $\varepsilon>0$, the set
\[
 S(\bar v,\varepsilon)=\{j\in\N:\bar v_j<\varepsilon\}
 \tag{4.1}
\]
is syndetic.
\item[(iii)] For every sequence $a=(a_m)_{m\in\N}$ of positive scalars and
every $\varepsilon>0$, the set
\[
 \left\{j\in\N:\inf_{m\in\N}a_m v^{(m)}_j<\varepsilon\right\}
 \tag{4.2}
\]
is syndetic.
\end{enumerate}
In particular, the criterion is independent of $p$.
\end{theorem}

\begin{proof}
\emph{(i)$\Rightarrow$(ii).}
Fix $\bar v\in\cV$ and $\varepsilon>0$.  Choose $c>0$ and put
\[
 \widetilde v=\bar v+c\,\mathbf1_{\{1\}}.
\]
Then $\widetilde v\in\cV$, because the added sequence has finite support,
and $\widetilde v_1>0$.  Consider the nonempty open sets
\[
 U=\{x:q_{\widetilde v}(x)<\varepsilon/2\},\qquad
 W=e_1+\{z:q_{\widetilde v}(z)<\widetilde v_1/2\}.
 \tag{4.3}
\]
If $n\in N_B(U,W)$, take $x\in U$ with $B^nx\in W$.  Since
$(B^nx)_1=x_{n+1}$, the second condition in (4.3) gives
\[
 |x_{n+1}-1|\widetilde v_1<\widetilde v_1/2,
 \quad\text{hence}\quad |x_{n+1}|>1/2.
\]
The first condition gives
\[
 |x_{n+1}|\widetilde v_{n+1}\leq q_{\widetilde v}(x)
 <\varepsilon/2.
\]
Consequently $\bar v_{n+1}\leq\widetilde v_{n+1}<\varepsilon$.  We have
proved
\[
 N_B(U,W)+1\subseteq S(\bar v,\varepsilon).
 \tag{4.4}
\]
The left side is syndetic by (i), and any superset of a syndetic set is
syndetic.  This proves (ii).  Notice that the finite-coordinate correction
also covers projective weights with $\bar v_1=0$.

\emph{(ii)$\Rightarrow$(i).}
Let $U,W$ be nonempty open subsets of $\kp$.  The finitely supported
sequences are dense, so there are finitely supported $x\in U$, $y\in W$ and
an absolutely convex zero neighborhood $G$ such that
\[
 x+G\subseteq U,\qquad y+G\subseteq W.
 \tag{4.5}
\]
Choose $s\geq1$ with both supports contained in $\{1,\ldots,s\}$, write
$y=\sum_{k=1}^s y_ke_k$, and set
\[
 M=1+\sum_{k=1}^s|y_k|,
 \qquad
 G_1=\bigcap_{r=0}^{s-1}B^{-r}(G/M).
 \tag{4.6}
\]
Continuity of $B$ makes $G_1$ a zero neighborhood.  By the projective
description, there are an integer $L$, weights
$\bar v^{(1)},\ldots,\bar v^{(L)}\in\cV$, and positive numbers
$\delta_1,\ldots,\delta_L$ such that
\[
 \bigcap_{\ell=1}^L\{z:q_{\bar v^{(\ell)}}(z)<\delta_\ell\}
 \subseteq G_1.
 \tag{4.7}
\]
Define a single projective weight by
\[
 \omega_j=
 \begin{cases}
 \left(\displaystyle\sum_{\ell=1}^L
 (\bar v^{(\ell)}_j/\delta_\ell)^p\right)^{1/p},&1\leq p<\infty,\\[3mm]
 \displaystyle\max_{1\leq\ell\leq L}
 \bar v^{(\ell)}_j/\delta_\ell,&p=0.
 \end{cases}
 \tag{4.8}
\]
Finite combination in (2.1) shows that $\omega\in\cV$.  Moreover,
$\omega_j<1$ implies $q_{\bar v^{(\ell)}}(e_j)<\delta_\ell$ for every
$\ell$, and hence $e_j\in G_1$ by (4.7).  Thus (ii) shows that
\[
 J=\{j:\omega_j<1\}
 \tag{4.9}
\]
is syndetic and $e_j\in G_1$ for every $j\in J$.

For $j\in J$ with $j\geq2s$, put
\[
 h_j=\sum_{k=1}^s y_ke_{j-s+k}
     =\sum_{k=1}^s y_kB^{s-k}e_j.
 \tag{4.10}
\]
By (4.6), every $B^{s-k}e_j$ belongs to $G/M$.  Absolute convexity and
$\sum_k|y_k|<M$ therefore give $h_j\in G$.  With $n=j-s\geq s$,
\[
 B^n(x+h_j)=B^nh_j=y.
 \tag{4.11}
\]
Equations (4.5) and (4.11) show that
\[
 \{j-s:j\in J,\ j\geq2s\}\subseteq N_B(U,W).
\]
Deleting a finite prefix and translating preserve syndeticity.  Hence the
right side is syndetic, proving (i).

\emph{(ii)$\Leftrightarrow$(iii).}
For every positive sequence $a$, the canonical weight in (3.1) belongs to
$\cV$, because for each fixed $r$,
\[
 \bar v^{,a}_j\leq a_rv^{(r)}_j\qquad(j\in\N).
\]
Thus (ii) implies (iii).  Conversely, if $\bar v\in\cV$, choose $a_m>0$
from (2.1).  Then
\[
 \bar v_j\leq\inf_m a_mv^{(m)}_j=\bar v^{,a}_j.
 \tag{4.12}
\]
The small-coordinate set in (4.2) is therefore contained in
$S(\bar v,\varepsilon)$.  Under (iii) the former is syndetic, so the latter
is syndetic as well.  This proves (ii).
\end{proof}

\section{Consistency checks and the remaining problem}

The theorem specializes correctly to the source's sufficient condition.
If for one row $m$ every set
$\{j:v^{(m)}_j<\eta\}$ is syndetic, then for
$\bar v\leq a_mv^{(m)}$ that row's $\varepsilon/a_m$-small set is contained
in $S(\bar v,\varepsilon)$.  It also specializes correctly to mixing: the
source characterizes mixing by $\bar v_j\to0$ for every $\bar v\in\cV$, in
which case every set in (4.1) is cofinite.  Unlike the fixed-row sufficient
condition, Theorem~\ref{thm:main} also accommodates the source's example
where mixing holds but no single row has the required syndetic small sets.

The source's Open Problem (2) remains open in this packet, but the theorem
gives an exact reduction.  Combining (iii) with the source's sequential
hypercyclicity criterion, one must construct a nuclear $V$ such that
\begin{enumerate}
\item for every positive $(a_m)$ and every $\varepsilon>0$, the set in
(4.2) is syndetic; while
\item for every $m$ and every thick $I\subseteq\N$, the sequence
$v^{(m)}_j$ fails to tend to zero as $j\to\infty$ through $I$.
\end{enumerate}
This separates the still-missing construction from the now-complete
ergodicity criterion.

\section{Verification and novelty note}

The proof was checked separately at the following points: projective weights
may vanish; the finite sum/max in (4.8) stays in $\cV$; the argument works
for both $p=0$ and finite $p$; and the shifts in (4.10)--(4.11) have the
correct indices.  Repository indexes and bounded exact-phrase,
title/author, arXiv, and projective-weight searches performed through
13 August 2026 did not locate a later answer to Open Problem (1).  This is a
novelty check, not a proof of priority.  Human review should focus first on
the reduction of an arbitrary zero neighborhood to (4.7) and on the
coordinate test (4.3).

\begin{thebibliography}{1}
\bibitem{BKP}
J.~Bonet, T.~Kalmes, and A.~Peris,
\emph{Dynamics of shift operators on non-metrizable sequence spaces},
arXiv:1911.07513v2; Rev. Mat. Iberoam. 37(6) (2021), 2373--2397.
\end{thebibliography}

\end{document}
